% ============================================================
% VI. DISCUSSION
% ============================================================

\section{Discussion}

The four exercises addressed in this report illustrate, in different ways, the same central idea: the digital representation of a signal — whether in the time domain (sampling), in amplitude (quantization), or in the execution order of an embedded program (concurrency) — always involves a potential loss of information or precision relative to the original continuous physical system, and that loss is controllable once its mechanism is understood.

In Item 1, the four distinct dominant frequencies obtained per segment (Table~\ref{tab:muestra01}) show that segmenting a non-stationary signal into short windows is a reasonable way to track how its spectral content changes over time, at the cost of frequency resolution within each window (shorter segment duration means lower frequency resolution, due to the inverse relationship between time and frequency resolution in the FFT).

In Item 2, the $f_s = 70$ Hz case experimentally confirms the aliasing phenomenon described in Section~\ref{sec:aliasing}: with fewer than two samples per cycle, the quantized samples do not unambiguously represent a 100 Hz sinusoid. This result, together with the aliasing-free $500$ and $1000$ Hz cases, demonstrates why the Nyquist theorem is not a theoretical footnote but a mandatory design condition in any data-acquisition system.

In Item 3, the most informative result was not that the three dominant frequencies differ (which was expected), but \emph{why} they differ: the violin and the cat produce peaks consistent with a perceptible fundamental tone, while the drum produces a peak dominated by the high-frequency content of its attack transient. This suggests that, while dominant frequency does allow these three sources to be \emph{differentiated} from one another, it should not be interpreted the same way for all of them: for tonal sources (violin, voices, most string or wind instruments) it approximates a perceptible fundamental note, whereas for broadband percussive sources (drum) it reflects the energy of the transient more than an identifiable tone. A more complete analysis should also consider the energy distribution across bands and the time envelope (attack, decay), not just the FFT's single peak.

In Item 4, all three sketches were flashed on a real ESP32-WROOM board and their serial output captured directly, so the results in Section~\ref{sec:embebidos} are genuine measurements rather than a model. The measured data confirms the qualitative structure expected from each mechanism's design: the busy loop's period scales linearly with $N$; the dual-core task keeps the nominal period essentially constant; and the interrupt keeps an exact timing reference while occasionally showing a processing backlog under a large $N$ — the single $N=10^7$ and $N=10^8$ backlog samples in \texttt{firmware/real\_logs/} are themselves evidence of this last effect, not noise to discard. A natural next step is to repeat the captures with a larger number of samples per run (more than the eight shown here) to obtain a tighter estimate of the jitter and backlog frequency for each mechanism, and to compare firmware 4(a)/4(b)/4(c) directly against an Arduino Uno implementation of the same busy-loop mechanism (item 4's numeral (a) and (c) for the Arduino Uno, which lacks a second core and has more limited interrupt/timer facilities than the ESP32).
